\documentclass[11pt]{article}
\usepackage[margin=1in]{geometry}
\usepackage{amsmath,amsthm,amssymb,booktabs}
\newtheorem{theorem}{Theorem}
\newtheorem{lemma}[theorem]{Lemma}
\newtheorem{corollary}[theorem]{Corollary}
\newtheorem{conjecture}[theorem]{Conjecture}
\theoremstyle{definition}
\newtheorem{definition}[theorem]{Definition}
\newtheorem{remark}[theorem]{Remark}
\newcommand{\E}{\mathbb{E}}
\newcommand{\dep}{\mathrm{dep}}
\title{How far can r-ary alias trees stretch an OSAM$^+$ trace?}
\author{Working note (SPARQ-OSAM), \today}
\date{}
\begin{document}
\maketitle

\begin{abstract}
We study the number of reads that the r-ary multi-write pointers of
\texttt{sam-model} (\texttt{RaryPointer}) add to an OSAM$^+$ trace. Against an
oblivious adversary, who fixes the access schedule in advance, and with a
randomized tree layout, we prove three things. The cost of objects with at most
$b$ aliases is deterministic. Each dereference deepens any alias by at most one
level. The total trace length exceeds its expectation by
$O\bigl(\sqrt{s\sum_o r_o^2}\bigr)$ except with probability $e^{-s}$, where the
ranges $r_o$ are computed from the schedule. For the \emph{binary} pointer we prove a
deterministic amortized bound, $2\log_2 f+1$ reads per dereference plus
$O(f)$ once, which holds even against adaptive adversaries. We state the stronger r-ary bound that the
experiments support ($O(c\sqrt{sK})$ for $K$ dereferences of high-fan-in
objects) as a theorem conditional on a bounded-influence lemma, and record the
measurements.
Section~\ref{sec:halving} gives an r-ary splay rule, path halving, with a
proven deterministic bound: $2\log_2 f+1$ reads per dereference, amortized.
It keeps every structural property of the current rule, and the experiments
support a $2\log_b f+O(1)$ bound, which we state as a conjecture.
\end{abstract}

\section{Model}

\begin{definition}[Alias tree]
An object with $f\ge 2$ aliases is a rooted tree whose root holds the value,
whose leaves are the aliases, and whose internal nodes have at most $b$ ordered
children ($b\ge 2$ even). A single pointer ($f=1$) is the root itself. The
\emph{install} layout is the balanced tree built by
\texttt{RaryPointer::install}, of height $H_0(f)=\lceil\log_b f\rceil$.
\end{definition}

\begin{definition}[Dereference]
Let alias $x$ be at depth $D$, on the path $u_0,\dots,u_D=x$ from the root
$u_0$. A dereference reads the $D+1$ cells $u_D,u_{D-1},\dots,u_0$, so its
cost is
\[ R = \dep(x)+1. \]
Let $A_j$ be the children of $u_j$ other than $u_{j+1}$, for $0\le j<D$. The
path is dissolved and rebuilt as follows:
\begin{enumerate}
\item $P \leftarrow (x', A_{D-1})$, where $x'$ is the fresh leaf of $x$.
\item For $j=D-2,D-3,\dots,0$: $P\leftarrow \mathrm{Pack}(P\circ A_j)$.
\item A fresh root receives the items of $P$ as its children.
\end{enumerate}
$\mathrm{Pack}(M)$ returns $M$ if $|M|\le b$. Otherwise it splits $M$ into
$m=\lceil |M|/b\rceil$ consecutive chunks of balanced sizes, and puts every chunk
of size $\ge 2$ under a fresh internal node.
\end{definition}

The standalone model in \texttt{src/aliasmodel.rs} implements exactly this
rule. It reproduces the library's read counts on 150{,}000 dereferences
($b\in\{2,\dots,30\}$, $f\le 3000$, three access orders;
\texttt{src/bin/validate\_model.rs}).

\begin{definition}[Randomized layout]
The initial assignment of aliases to install leaves is a uniform bijection.
Every $\mathrm{Pack}$, and the final placement under the root, first applies a
fresh uniform permutation to its items. All other randomness is independent
across dereferences.
\end{definition}

\begin{definition}[Oblivious adversary, trace length]
The adversary fixes a schedule $\sigma=((o_1,a_1),\dots,(o_K,a_K))$ of
(object, alias) dereferences before any randomness is drawn. It may depend on
public parameters ($b$ and every fan-in $f_o$) but not on the layout. The trace
length is $T(\sigma)=\sum_{t} R_t$.
\end{definition}

\begin{remark}
In the SAM model, the read count is a deterministic function of the schedule
and the layout. Path OSAM$^+$'s own randomness (leaf positions, eviction)
never changes it. With the library's deterministic layout, $T(\sigma)$ is
therefore a fixed number for each $\sigma$. "Exceeding the expectation" is
only meaningful once the layout is randomized, which is why the randomized
layout is part of the model.
\end{remark}

\section{Structural lemmas (unconditional)}

\begin{lemma}[At most $b$ items pending]\label{lem:pending}
During a rebuild, $|P|\le b$ after step 1 and after every $\mathrm{Pack}$. In
particular, the root never needs extra packing.
\end{lemma}
\begin{proof}
After step 1, $|P|\le 1+(b-1)=b$. If $|P|\le b$, then $|P\circ A_j|\le 2b-1$.
$\mathrm{Pack}$ then either returns at most $b$ items unchanged, or returns
$\lceil (2b-1)/b\rceil = 2\le b$ chunks.
\end{proof}

\begin{lemma}[Unit depth growth]\label{lem:growth}
Consider one dereference of alias $x$ in an object. Every other alias $y$ of
that object satisfies $\dep'(y)\le \dep(y)+1$, and $x$ itself satisfies
$\dep'(x)\le\dep(x)$. Aliases of other objects are unaffected. This holds for
every layout, deterministic or randomized.
\end{lemma}
\begin{proof}
Every leaf other than $x$ lies in exactly one side subtree $A_j$, whose
internal shape is untouched. Only its root moves. An item's new depth is $1$
(child of the new root) plus the number of $\mathrm{Pack}$s that wrapped it
under a fresh node. Each $\mathrm{Pack}$ wraps an item at most once.

For $j\le D-2$, $A_j$ enters $P$ at level $j$, so it takes part in the
$\mathrm{Pack}$s at levels $j,j-1,\dots,0$. That is $j+1$ of them, so its
root's new depth is at most $j+2$. The old depth was $j+1$.

$x'$ and $A_{D-1}$ take part in all $D-1$ $\mathrm{Pack}$s, so their new depth
is at most $D$. The old depth was $D$.
\end{proof}

\begin{corollary}[Flat regime]\label{cor:flat}
If $f\le b$, every dereference costs exactly $2$ reads, forever (exactly $1$
read if $f=1$).
\end{corollary}
\begin{proof}
The install tree is flat, so $D=1$. Then $P=(x',A_0)$ holds all $f\le b$
leaves, no $\mathrm{Pack}$ occurs, and the new root again has all $f$ leaves as
children. Induct on the number of dereferences.
\end{proof}

\begin{corollary}[Schedule-computable range]\label{cor:range}
Fix $\sigma$ and an object $o$ with $f_o>b$. Let $t$ be a dereference of alias
$a$ of $o$. Let $g_t$ be the number of dereferences of $o$ since $a$'s previous
dereference $t'$, or since the start if there is none. Then
\[
2\le R_t\le U_t(\sigma):=\begin{cases}U_{t'}(\sigma)+g_t & \text{if } t' \text{ exists},\\
H_0(f_o)+1+g_t & \text{otherwise.}\end{cases}
\]
\end{corollary}
\begin{proof}
Apply Lemma~\ref{lem:growth} once per intervening dereference of $o$, starting
from the install height. The lower bound is $\dep\ge 1$ for $f\ge 2$.
\end{proof}

\section{Concentration}\label{sec:concentration}

\begin{theorem}[Concentration against oblivious adversaries]\label{thm:main}
Assume a randomized layout, and fix any schedule $\sigma$. Let $O_{>}$ be the
set of objects with $f_o>b$. Let $n_o$ be the number of dereferences of $o$ in
$\sigma$, and $r_o=\sum_{t\,:\,o_t=o}(U_t(\sigma)-2)$. Then for every
$\lambda>0$,
\[
\Pr\bigl[T(\sigma)\ \ge\ \E[T(\sigma)]+\lambda\bigr]\ \le\
\exp\!\Bigl(-\frac{2\lambda^2}{\sum_{o\in O_>} r_o^2}\Bigr).
\]
Equivalently, with probability at least $1-e^{-s}$,
$T(\sigma)\le \E[T(\sigma)]+\sqrt{\tfrac{s}{2}\sum_{o\in O_>} r_o^2}$.
\end{theorem}
\begin{proof}
Write $T(\sigma)=C(\sigma)+\sum_{o\in O_>}T_o$. The term $C(\sigma)$ collects
the objects with $f_o\le b$, and is deterministic by
Corollary~\ref{cor:flat}. $T_o$ depends only on $\sigma|_o$, on the initial
bijection of $o$, and on the permutations drawn during $o$'s own rebuilds. The
schedule is fixed in advance, and fresh randomness is drawn per rebuild, so
these random inputs are disjoint across objects. The $T_o$ are therefore
mutually independent. (Drawing everything from one shared stream gives the
same joint law, because each draw uses fresh bits.) By
Corollary~\ref{cor:range}, $T_o\in[2n_o,\,2n_o+r_o]$. Hoeffding's inequality
for independent bounded summands finishes the proof.
\end{proof}

\paragraph{Reading the bound.}
Suppose each object with $f>b$ is dereferenced at most once in the run
($r_o=H_0-1$) over $K$ such dereferences. The excess is then at most
$(H_0-1)\sqrt{sK/2}$, i.e.\ it grows like $\sqrt K$. With $b=30$, $f=952$
($H_0=3$), $K=10^4$ and $s=40$ (failure probability $2^{-57}$), the excess is at
most 894 reads, on an expected total of about $4.5\cdot 10^4$ (2\%). The bound
degrades for \emph{hot} objects: an object dereferenced $n$ times with long gaps
has $r_o=\Theta(n^2)$ in the worst case. That is where
Theorem~\ref{thm:cond} is needed.

\begin{definition}[Influence]
Let $\rho_0$ be the initial layout and $\rho_t$ the randomness of the $t$-th
dereference. The influence of step $t$ is
\[
c_t=\sup_{\rho_{<t}}\ \sup_{\rho_t,\rho_t'}\
\Bigl|\E[T\mid \rho_{<t},\rho_t]-\E[T\mid \rho_{<t},\rho_t']\Bigr|.
\]
\end{definition}

\begin{theorem}[Conditional $\sqrt K$ bound]\label{thm:cond}
If $c_t\le c$ for every $t$, then for every schedule $\sigma$ with $K$
dereferences of objects in $O_>$,
\[
\Pr\bigl[T(\sigma)\ge \E[T(\sigma)]+\lambda\bigr]\le \exp\!\Bigl(-\frac{\lambda^2}{2Kc^2}\Bigr).
\]
\end{theorem}
\begin{proof}
$M_t=\E[T\mid\rho_0,\dots,\rho_t]$ is a Doob martingale in which only the
$K$ steps that touch $O_>$ carry randomness. Its increments are bounded by
$c_t\le c$. Apply Azuma--Hoeffding.
\end{proof}

\begin{conjecture}\label{conj}
For the randomized layout, $c=O(b\,H_0(f))$, independent of the schedule
and of $K$. Also, for every oblivious schedule,
$\E[R_t]\le H_0(f)+1+\gamma_b$ for a constant $\gamma_b$.
\end{conjecture}

The second half of the conjecture is what padding needs. The public budget has
to cover $\max_\sigma \E[T(\sigma)]$ plus the slack.

\section{Evidence}

\paragraph{Concentration} (\texttt{src/bin/concentration.rs}). This uses 100
independent layouts per fixed schedule and a window of $k$ dereferences of one
object after a warm-up of $2f$ dereferences. Across all $(b,f)\in\{(4,300),(16,952),(30,952),(30,3000)\}$ and all six schedules (cyclic,
stride, uniform, Zipf, bursts, hot-set-then-scan), we observe:
\begin{itemize}
\item $\mathrm{sd}(T_k)/\sqrt k\in[0.59,1.32]$ for every schedule except the phased hot-set-then-scan one ($0.81$--$4.13$).
\item The maximum over 100 layouts is at most $\E+3.1\,\mathrm{sd}$.
\end{itemize}
This is the $\sqrt K$ scaling of Theorem~\ref{thm:cond} with $c\approx 1$, even
though a single object is dereferenced $10^4$ times, far outside the regime
where Theorem~\ref{thm:main} is sharp.

\paragraph{Influence} (\texttt{src/bin/influence2.rs}). We branch at a random
step, then use common random numbers for 400 futures. The signed influence
estimates for $b=30$, $f=952$ at look-ahead $L=10,100,1000,4000$ are
below one read in absolute value at $L\le 100$. At $L=1000$ and $4000$ they
are at most $9$ reads for the uniform and cyclic schedules, all within
$2.3$ standard errors of $0$. They are larger but still noisy ($\le 22$
reads) for hot-set-then-scan. Under this coupling the two copies do not
re-synchronize, so the per-path differences grow like $\sqrt L$ while their
means stay small. A proof of Conjecture~\ref{conj} needs a coupling that
merges two layouts once they have pushed the same set of items through a
$\mathrm{Pack}$.

\paragraph{Expectation} Mean reads per dereference with the randomized layout
(deterministic library in parentheses):

\medskip
\begin{tabular}{lccccc}
\toprule
$(b,f)$ & $H_0+1$ & uniform & cyclic & Zipf & bursts\\
\midrule
$(4,300)$ & 6 & 8.11 (7.70) & 8.33 (6.43) & 7.04 (6.80) & 2.87 (2.20)\\
$(16,952)$ & 4 & 5.51 (5.07) & 5.56 (4.79) & 4.92 (4.66) & 2.44 (2.04)\\
$(30,952)$ & 4 & 4.79 (4.48) & 4.82 (4.05) & 4.40 (4.18) & 2.34 (2.02)\\
$(30,3000)$ & 4 & 5.55 (5.07) & 5.58 (3.59) & 4.95 (4.60) & 2.45 (2.03)\\
\bottomrule
\end{tabular}
\medskip

\noindent Randomizing the layout costs 5--55\% in expectation, but it turns
worst-case-over-schedules padding into expectation-plus-$O(\sqrt K)$ padding.

\paragraph{The deterministic potential fails} (\texttt{src/bin/potential.rs}).
With the library's deterministic layout, the rank potential
$\Phi=c\sum_v\log_2 s(v)$ does not certify an amortized bound of the form
$R_i+\Phi_i-\Phi_{i-1}\le O(\log f)$ for any tested scale $c$. Single rebuilds
can raise $\Phi$ by 12--25 bits, and an adversary that always picks the deepest
alias is not the worst case; cyclic scans are. This is one more reason to use
the randomized layout together with an oblivious adversary.


\section{The binary pointer: a deterministic bound}

The binary multi-write pointer (\texttt{MultiWritePointer}) admits a much
simpler and stronger statement. Its bound is deterministic, so it needs no
randomized layout and holds even against \emph{adaptive} adversaries.

\begin{definition}[Binary alias tree and rebuild]
The tree is full binary: every cell stores its parent and its sibling, and
leaves are aliases. The install layout splits $f$ aliases into
$\lfloor f/2\rfloor$ and $\lceil f/2\rceil$ recursively. Dereferencing alias
$x$ maintains a state $(X,Y,P)$ of two new subtrees and one old node with
$\mathrm{leaves}(X)\cup\mathrm{leaves}(Y)=\mathrm{leaves}(P)$. It starts from
$(\{x'\},\ \mathrm{sib}(x),\ \mathrm{par}(x))$ after reading $x$. Each
iteration reads $P$:
\begin{itemize}
\item If $P$ is the root, a new root receives children $X$ and $Y$. Stop.
\item Otherwise read $G=\mathrm{par}(P)$, and let $S_P=\mathrm{sib}(P)$.
  \begin{itemize}
  \item If $G$ is the root, form $N=(Y,S_P)$; a new root receives $X$ and $N$. Stop.
  \item Otherwise let $S_G=\mathrm{sib}(G)$, form $L=(X,Y)$ and $R=(S_P,S_G)$, and continue with $(L,R,\mathrm{par}(G))$.
  \end{itemize}
\end{itemize}
Every path cell is read once, so the cost is again $\dep(x)+1$. This is
path halving: two levels are consumed per iteration and one is created.
\end{definition}

\begin{lemma}[Access lemma]\label{lem:access}
Let $\Phi=\sum_{v\text{ internal}}\log_2 s(v)$, where $s(v)$ is the number of
aliases below $v$. One dereference in an object with $f$ aliases satisfies
\[ R\ \le\ 2\log_2 f+1+\Phi-\Phi'. \]
\end{lemma}
\begin{proof}
Let $\mu=\log_2 s(P)$ for the current state, so $\mu\ge 1$ after the first
read.

In a middle iteration, write $a=s(P)$, $b=|S_P|$ and $c=|S_G|$. The node $L$
has exactly $P$'s leaves, and $R$ replaces $G$, so
$\Delta\Phi=\log_2(b+c)-\log_2(a+b)$. By AM--GM,
$a(b+c)\le\bigl((a+b+c)/2\bigr)^2$, so
\[
2+\Delta\Phi\ \le\ 2\log_2(a+b+c)-\log_2 a-\log_2(a+b)\ \le\ 2(\mu'-\mu),
\]
where $\mu'=\log_2(a+b+c)=\log_2 s(\mathrm{par}(G))$.

The root case costs one read, with $\Delta\Phi=0$ and $\mu=\log_2 f$. The
$G$-root case costs two reads, with
$\Delta\Phi=\log_2(|Y|+|S_P|)-\mu\le\log_2 f-\mu$.

Telescoping from $\mu_0\ge 1$, the total is at most $2\log_2 f$ in the first
case and $1+\mu+\log_2 f\le 2\log_2 f+1$ in the second.
\end{proof}

\begin{theorem}[Deterministic trace bound, binary pointers]\label{thm:binary}
For every sequence of $k$ dereferences of one object, chosen adaptively or
not, starting from its install layout,
\[
T\ \le\ k\,(2\log_2 f+1)+\Delta(f),\qquad \Delta(f)=B(f)-(f-1)\le 2f,
\]
where $B(1)=0$ and $B(f)=\log_2 f+B(\lfloor f/2\rfloor)+B(\lceil f/2\rceil)$
is the install potential. Over many objects,
$T(\sigma)\le\sum_t(2\log_2 f_{o_t}+1)+\sum_{o\text{ touched}}\Delta(f_o)$. A
window that starts in an arbitrary state instead of the install layout pays at
most $(f-1)(\log_2 f-1)$ instead of $\Delta(f)$.
\end{theorem}
\begin{proof}
Sum Lemma~\ref{lem:access}; the potentials telescope. Every internal node has
$s\ge 2$ and there are $f-1$ of them, so $\Phi\ge f-1$. Also
$\Phi\le(f-1)\log_2 f$. Induction gives $B(f)\le 3f-\log_2 f-3$, since
$2\log_2 f-\log_2(\lfloor f/2\rfloor\lceil f/2\rceil)\le 2+\log_2\tfrac43<3$.
Hence $\Delta(f)\le 2f$. Numerically, $\Delta(f)\le 1.06f$ for all
$f\le 2\cdot10^5$.
\end{proof}

\paragraph{Check} (\texttt{src/bin/potential\_binary.rs}). Lemma~\ref{lem:access}
held on all 126{,}000 dereferences of the library's
\texttt{MultiWritePointer}: $f\in\{2,\dots,3000\}$, with random, cyclic,
reverse, Zipf, bursty, and always-deepest (adaptive) schedules. Every step had
at least one read of slack.

\paragraph{Why the r-ary rule resists this argument.} The proof uses two
properties of path halving: the new left node inherits its old parent's rank
exactly, and only one new node absorbs the side subtrees of two levels. The
r-ary rebuild carries up to $b$ pending items per level and re-chunks them,
so it does neither. We measured single rebuilds raising $\Phi$ by 12--25 bits.

\paragraph{Price.} Binary trees cost $\approx\log_2 f$ reads per dereference,
against $\approx\log_b f$. For $f=952$ the binary pointer's mean is 12.9
(uniform), 6.7 (cyclic) and 2.3 (bursts) reads, and Theorem~\ref{thm:binary}
certifies 20.8 per dereference. The r-ary pointer (b = 30) averages 4.5 on the
uniform schedule, and a shape-preserving r-ary rebuild costs exactly 4.

\section{What the server sees}
\label{sec:rw}

Standard OSAM$^+$ (the setting of the Concrete OSAM paper) performs every
write by reading a dummy path, so it leaks the number of reads \emph{plus}
writes. There the binary pointer costs exactly $C=3R-1$ per dereference. The
r-ary backend used here (\texttt{MULTI\_WRITE\_RARY}) is different. Writes are
local: they go to the stash and are drained by the $P$ eviction paths that each
read downloads, with $P$ much smaller than $b$ (we use $P=1$). The server sees
only reads, each of $1+P$ paths, plus public flushes during bulk builds. So
$R$ is the leaked quantity, and a large $b$ saves both round trips and
bandwidth. All the rules in this note write $O(b)$ cells per read, which is
the property the backend relies on to keep the stash bounded. The alias model
now also counts writes. It matches the library write for write on the same
150{,}000 dereferences (\texttt{validate\_model}).

\section{An r-ary splay rule with a proven bound: path halving}
\label{sec:halving}

The current rule dissolves the whole path and re-chunks it, and no rank
potential certifies it (Section~5). We now keep the splay behavior but
restructure like the binary pointer: two path nodes per step, halving the
path. The difference from the binary rule is that the carried set may hold up
to $b$ items.

\paragraph{The rule.} Let $s(v)$ be the number of aliases below $v$. Write
$\mathrm{sib}(v)$ for the other children of $v$'s parent. A dereference of
alias $a$ reads $a$'s cell and replaces $a$ by a fresh leaf. It then keeps a
\emph{carried set} $K$: the new child set of the current path node $p$,
initially the children of $a$'s parent with $a$ replaced. Repeat:
\begin{enumerate}
\item Read $p$. If $p$ is the root, write a new root with children $K$. Stop.
\item Let $S_1=\mathrm{sib}(p)$, and read the parent $p'$ of $p$. If $p'$ is
  the root, write a new root with children $K\cup S_1$ if $|K|+|S_1|\le b$,
  and otherwise with children $\{\ell\}\cup S_1$, where $\ell$ is a new node
  over $K$. Stop.
\item Let $S_2=\mathrm{sib}(p')$ and $p''$ the parent of $p'$. The new child
  set $K'$ of $p''$ is:
  \begin{itemize}
  \item[(A0)] $K\cup S_1\cup S_2$, if it has at most $b$ items;
  \item[(A)] otherwise $\{\ell\}\cup S_1\cup S_2$, with $\ell=\mathrm{node}(K)$,
    if that has at most $b$ items;
  \item[(B)] otherwise $\{\ell,r\}\cup(S_2\setminus D)$. Here $D$ holds the
    $\max(0,|S_2|+2-b)\le 1$ smallest members of $S_2$, and
    $r=\mathrm{node}(S_1\cup D)$, or $r$ is that single member when
    $|S_1\cup D|=1$.
  \end{itemize}
  Discard $p$ and $p'$, set $p\leftarrow p''$ and $K\leftarrow K'$, and
  repeat.
\end{enumerate}

\begin{theorem}[r-ary path halving]\label{thm:halving}
Let $\Phi=\sum_{v\ \mathrm{internal}}\log_2 s(v)$ and $f\ge2$. Every
dereference satisfies:
\begin{enumerate}
\item (single read, bounded groups) every cell on the path is read once, the
  carried set and every new group have at most $b$ members, and every node
  keeps at least 2 children;
\item (unit depth growth) every other alias deepens by at most one level, and
  the accessed alias never deepens;
\item (flat regime) if $f\le b$, it performs exactly 2 reads and the tree stays
  flat;
\item (access lemma) it performs $R\le 2\log_2 f+1+\Phi-\Phi'$ reads.
\end{enumerate}
Consequently, for every $k$ and every sequence of $k$ dereferences of an
object installed as $T_0$, even one chosen adaptively,
\[
\Pr\Big[\sum_t R_t > k(2\log_2 f+1)+\Phi(T_0)-\tfrac{f-1}{b-1}\Big]=0 .
\]
\end{theorem}
\begin{proof}
(1) holds by construction. $|K|\le b$ holds initially, and (A0), (A) and (B)
each produce at most $b$ items. In (B), $|K'|=2+|S_2|-|D|\le b$ and
$|S_1\cup D|\le (b-1)+1$. Every new node gets at least 2 children.

(2) Let $d$ be the depth of $p$. A step moves $K$'s members from depth $d+1$
to depth $d$ or less. It moves $S_1$ from depth $d$ to $d$ (inside $r$) or to
$d-1$. It moves $D$ from $d-1$ to $d$, and it leaves $S_2\setminus D$ in
place. Everything a step places in $K'$ only moves up afterwards, and the
final steps do not deepen anything. Every off-path subtree occurs in exactly
one $S_1$ or $S_2$, so it deepens at most once.

(3) If $f\le b$, the root's children are all aliases, so step 1 applies at
once: 2 reads, and the tree is again flat.

(4) Let $x=s(p)$ and $\mu=\log_2 x$. Initially $\mu\ge1$. $K$ covers the same
aliases as the old $p$, and $\Delta\Phi$ counts the removed and the created
nodes. Step 1 performs 1 read, and $\Delta\Phi=0$. Step 2 performs 2 reads,
and $\Delta\Phi\le0$. In step 3, let $y=s(S_1)$, let $\delta=s(D)$, and let
$\mu'=\log_2 s(p'')$, where $s(p'')\ge x+y+\delta$. The step performs 2 reads.
\begin{itemize}
\item (A0): $\Delta\Phi=-\log_2x-\log_2(x+y)\le -2$, so $2+\Delta\Phi\le0$.
\item (A): $\Delta\Phi=-\log_2(x+y)$. Every node has at least 2 children, so
  $x\ge2$, $y\ge1$ and $s(p'')\ge4$. Hence $2+\Delta\Phi\le2(\mu'-\mu)$ (check
  $x+y=3$ separately).
\item (B): $\Delta\Phi=\log_2(y+\delta)-\log_2(x+y)$. By AM--GM,
  $x(y+\delta)\le((x+y+\delta)/2)^2$, so
  \[
  2+\Delta\Phi\le 2\log_2(x+y+\delta)-\log_2x-\log_2(x+y)\le2(\mu'-\mu).
  \]
  If $r$ is a single member, $\Delta\Phi=-\log_2(x+y)$, as in (A).
\end{itemize}
Summing, the steps telescope from $\mu\ge1$ to at most $\log_2 f$. With the
leaf read and the final step, $R+\Delta\Phi\le 1+2(\log_2f-1)+2=2\log_2 f+1$.
For the corollary, sum the lemma, and use that at the end every internal node
has size at least 2 and there are at least $(f-1)/(b-1)$ of them.
\end{proof}

The bound is deterministic and needs no randomness or assumption on the
adversary. It is the r-ary counterpart of the binary theorem
(Theorem~\ref{thm:binary}), and it keeps every structural property of the
current rule. For the installed tree, $\Phi(T_0)=O(f\log_2 b/b)$ is small.

We checked all four claims on 180{,}000 dereferences (\texttt{halving}): $b\in\{4,8,30\}$,
$f\in\{16,100,952,3000\}$, and uniform, cyclic, Zipf, bursty and
deepest-alias schedules. The access lemma always held, with slack at least
6 reads. The largest depth growth of another alias was 1, and $f\le b$ always
cost 2.

\paragraph{The proven bound is in base 2; the data say base $b$.} The telescoping
above gains only a factor 2 in size per two reads, as in the binary pointer.
With the potential $2\sum\log_b s(v)$ instead, the largest amortized cost of a
single dereference, $R+2\Delta\Phi_b$, was:

\begin{center}\small
\begin{tabular}{rr|c|cc|c}
$b$&$f$&$2\log_b f$&fixed schedules (6 kinds)&greedy adaptive adversary&$2\log_2f+1$\\\hline
4&952&9.9&8.6&8.7&20.8\\
4&3000&11.6&10.4&11.1&24.1\\
8&952&6.6&7.0&6.7&20.8\\
30&952&4.0&5.0&5.0&20.8\\
30&3000&4.7&5.0&5.0&24.1
\end{tabular}
\end{center}
The greedy adversary (\texttt{halving\_adv}) tries the deepest alias and 40
random ones at every step, on a clone of the tree, and dereferences the one
with the largest amortized cost.

\begin{conjecture}\label{conj:halving}
For r-ary path halving, $R\le 2\log_b f+c+2(\Phi_b-\Phi_b')$ for an absolute
constant $c$ (the data suggest $c\le1$), where $\Phi_b=\sum_v\log_b s(v)$.
\end{conjecture}

\emph{Where a proof must work.} In a type-(B) step, write $u=y+\delta$ for
the size of $r$ and $z$ for the total size of $S_2$. The base-$b$ step
inequality is $s(p'')(x+y)/(x\,u)\ge b$. It fails in a binary-like
configuration: $p$ has one sibling subtree of about its own size ($y\approx
x$), while $p'$ has many small siblings ($z$ small, but $|S_1|+|S_2|\ge b$).
Then the step doubles the size for two reads. Such steps cost
$2-4\log_b2$ each, beyond the base-$b$ bound. A proof needs a potential that
also charges shapes where a node has one child much larger than the rest.
Neither the fixed schedules nor the greedy adversary produced a run of such
steps.

\paragraph{Against the current rule.} Reads per dereference (mean / max,
3{,}000 steps):
\begin{center}\small
\begin{tabular}{rr|cc|cc|cc|cc}
&&\multicolumn{2}{c|}{uniform}&\multicolumn{2}{c|}{cyclic}&\multicolumn{2}{c|}{deepest}&\multicolumn{2}{c}{bursty}\\
$b$&$f$&halving&current&halving&current&halving&current&halving&current\\\hline
4&952&6.5/9&9.0/16&4.3/13&5.4/45&7.0/8&9.8/11&3.3/11&3.1/25\\
8&3000&5.3/7&7.2/12&3.5/6&3.2/35&6.0/7&7.7/9&3.2/7&2.8/17\\
30&952&3.1/5&4.4/8&3.1/5&2.7/10&4.0/5&4.7/6&3.0/5&2.3/11\\
30&3000&4.0/5&5.0/9&3.3/5&2.5/7&5.0/5&5.6/7&3.1/5&2.4/10
\end{tabular}
\end{center}
Halving has a lower mean on uniform and deepest-alias schedules, and a much
lower maximum everywhere: its worst single dereference stays near
$\log_b f+2$. It is slower on bursts, because a re-accessed alias only moves
up half the path, so it costs 3 reads rather than 2. A final step that
collapses into the root whenever the carried set fits would recover part of
that.

\section{A balanced r-ary pointer with downward pointers}
\label{sec:balanced}

For static structures, where queries only dereference, and for any setting
that needs a worst-case bound, we give up splaying and keep each alias tree
\emph{complete}. The shape is a function of the alias count $f$ alone: all
leaves are at depth $H=\lceil\log_b f\rceil$ ($H=0$ for $f=1$, where the alias
is the root), and leaf $j$ follows the base-$b$ digits of $j$
(\texttt{src/bptr.rs}).

\paragraph{Cells.} A root holds the value, $f$, $H$ and the address of its
child list. Every alias and every internal node has a cell holding its
parent's address and the parent's whole child group, as pairs
$(\text{node},\text{down})$. An internal node $v$ also has a child-list cell
$\mathrm{Kids}(v)$ that holds the same group. This is the downward pointer.
Because a member's cell is a function of its parent and its group, a
dereference can rewrite every off-path member without reading it. Child lists
are only rewritten during dereferences, never read, so their addresses stay
put until a copy or delete reads them.

\begin{theorem}[Balanced pointer]\label{thm:balanced}
Starting from a bulk install or from $\mathrm{New}$, and after any sequence
of dereferences, copies and deletes, even one chosen adaptively, the alias
tree is the complete tree for the current $f$, and:
\begin{enumerate}
\item a dereference performs exactly $H+1$ reads and at most $(b+1)H+1$
  writes (at most $b+1$ per read). It performs fewer reads only when the walk
  meets a root that is open in the client cache;
\item a copy and a delete each perform exactly $2H$ reads, or 1 when $H=0$;
\item every address is read at most once, and an alias's cell is read only by
  its own holder.
\end{enumerate}
Hence, for every $k$, $\Pr\big[\sum_t R_t>\sum_t(\lceil\log_b
f_t\rceil+1)\big]=0$ over any $k$ dereferences.
\end{theorem}
\begin{proof}[Proof sketch]
\emph{Dereference.} The walk reads the alias and each ancestor once. Each read
cell gets a fresh address. For each path node $v$, the group of $v$'s
children is known, because it is stored in the cell of the child we came
from. So all of $v$'s children are rewritten with $v$'s new address, and
$\mathrm{Kids}(v)$ is rewritten in place. No node moves, so the shape is
unchanged.

\emph{Copy.} After the upward walk, the new alias goes to position $f$. The
walk goes down from the root along the digits of $f$. It reads
$\mathrm{Kids}(u)$ for each frontier node $u$ that is not on the alias's path
(at most $H-1$ of them), gives that child list a fresh address, and rewrites
$u$'s group. Missing frontier nodes are created. If $f=b^H$, a new level is
inserted below the root.

\emph{Delete.} The upward walk, then the downward walk to position $f-1$
(again at most $H-1$ child-list reads). The last alias is moved into the
hole: its cell is rewritten in place, so its holder's address stays valid.
Emptied frontier nodes are removed. If the root is left with a single
internal child, the tree is contracted, which needs at most one more
child-list read, and only when the frontier walk needed none. When one alias
remains, its cell becomes the root.

Copies and deletes are padded with dummy reads to exactly $2H$. Every read in
these walks is of a path node or a child list, never of another holder's
alias.
\end{proof}

The implementation passes randomized tests of every claim, with exact read
counts: dereferences for $b=2,\dots,30$ and $f\le 952$; 3{,}000 random
copies, deletes, puts and gets per fanout; copies that grow the tree while its
root is cached; and mixed workloads through sam-model's client cache. With
$b=30$ and 256-byte blocks (a node cell is $6+8b$ bytes), SPARQ-OSAM's
per-query budgets now hold with raw fan-in. On seven datasets, with 500
queries per tail, all answers were correct and no query exceeded its
budget.

Compared with the splay rules of Sections~5--7, the balanced pointer gives up
adaptivity and zero-read copies. In return it has exact worst-case costs and
no trace variance from the pointer state.

\end{document}
